\documentclass[11pt]{article}
\usepackage[margin=1in]{geometry}
\usepackage{amsmath,amssymb,amsthm}
\usepackage[hidelinks]{hyperref}
\newtheorem{theorem}{Theorem}
\newtheorem{lemma}{Lemma}
\DeclareMathOperator{\exc}{exc}
\DeclareMathOperator{\anxc}{anxc}
\DeclareMathOperator{\fix}{fix}
\DeclareMathOperator{\sgn}{sgn}
\newcommand{\A}{\mathcal A_n^S}
\title{Excedance--antiexcedance elements commute\\
\large MathDB 369524}
\author{Principia Math in Progress}
\date{12 September 2026}
\begin{document}
\maketitle

\paragraph{Task 1. Solve the original conjecture.}
For \(w\in S_n\), write
\[
\exc(w)=\#\{i:w(i)>i\},\qquad
\anxc(w)=\#\{i:w(i)<i\}.
\]
For a commutative ring \(\mathbf k\), let
\[
X_{a,b}=\sum_{\substack{w\in S_n\\\exc(w)=a,\ \anxc(w)=b}}w
\quad\text{in }\mathbf k[S_n].
\]
Grinberg's Conjecture~3.10.22~\cite{grinberg} asks whether all these
elements commute for fixed \(n\). The answer is yes.

\begin{theorem}
For every \(n\), every commutative coefficient ring \(\mathbf k\), and all
nonnegative \(a,b,c,d\),
\[
X_{a,b}X_{c,d}=X_{c,d}X_{a,b}.
\]
Over \(\mathbb C\), the algebra generated by the \(X_{a,b}\) is the image
of the homogeneous Bethe algebra \(\A\) under the sign automorphism
\(\varepsilon(w)=\sgn(w)w\).
\end{theorem}

\paragraph{The known Bethe algebra.}
We use the following published results of Mukhin, Tarasov and
Varchenko~\cite{mtv}, with numbering from their arXiv v4.
Specialize \(z_1=\cdots=z_n=0\) and \(\hbar=1\).
For \(R\subseteq[n]\), put
\[
e_R=\sum_{\sigma\in S_R}\sgn(\sigma)\sigma,
\]
where permutations of \(R\) fix its complement. Their formula (6.7) is
\begin{equation}\label{eq:transfer}
S_k(u)=
\sum_{\substack{R\subseteq[n]\\|R|=k}} e_R
\prod_{a\notin R}^{\longleftarrow}
\left(u+\sum_{\substack{r\in R\\r<a}}(r,a)\right).
\end{equation}
The arrow means decreasing \(a\) from left to right, and multiplication
of permutations means composition with the rightmost factor acting first.
The convention \(S_0(u)=u^n1\) agrees with this formula.
Their Proposition~6.8 says that the coefficients of these polynomials
generate \(\A\); Theorem~6.18 says it is commutative.
Section~8 identifies this specialization as the homogeneous Bethe algebra.
These established facts are the only non-elementary input.

Set
\[
\widehat S_n(u,x)=\sum_{k=0}^nS_k(u)x^k,
\]
as in their formula (6.9). Every coefficient lies in \(\A\), and
every monomial \(x^ku^m\) occurring has \(k+m\leq n\).

\begin{lemma}\label{lem:coefficient}
For every \(w\in S_n\),
\begin{equation}\label{eq:coefficient}
[w]\widehat S_n(u,x)=
\sgn(w)x^{\anxc(w)}(x-1)^{\exc(w)}(u+x)^{\fix(w)}.
\end{equation}
\end{lemma}
\begin{proof}
Partition the \emph{values} of \(w\) into
\[
D=\{w(i):w(i)<i\},\quad E=\{w(i):w(i)>i\},\quad F=\{i:w(i)=i\}.
\]
Write \(b=|D|\), \(a=|E|\), \(f=|F|\).
We first prove that terms in \eqref{eq:transfer} producing \(w\) are in
bijection with root sets
\begin{equation}\label{eq:roots}
R=D\ \dot\cup\ J\ \dot\cup\ H,\qquad J\subseteq E,\quad H\subseteq F.
\end{equation}

Start with a permutation \(\sigma\) of \(R\). As the factors in
\eqref{eq:transfer} are multiplied on its right in decreasing \(a\),
choosing the transposition \((r,a)\) inserts the currently fixed value
\(a\) immediately after \(r\) in its cycle. Choosing \(u\) leaves \(a\)
fixed. For each retained \(r\), the inserted values immediately following
it therefore occur in increasing order, and all exceed \(r\).
Every nonfixed value outside \(R\) has a smaller predecessor in the final
cycle, so it lies in \(E\). This proves the necessity of \eqref{eq:roots}.

Conversely, suppose \(D\subseteq R\). Every nontrivial cycle of \(w\)
contains a value of \(D\), since its least value has a larger predecessor.
Cut such a cycle at its retained values. Each intervening run of omitted
values is strictly increasing: an omitted value has a smaller predecessor,
as it belongs to \(E\). Delete these runs to obtain the unique permutation
\(\sigma\) of \(R\). For each omitted nonfixed value \(a\), choose
\((r,a)\), where \(r\) is the retained value just before its run.
Here \(r<a\), so the factor is allowed. Multiplying these factors in
decreasing order reconstructs every run, hence reconstructs \(w\).
Omitted fixed points choose \(u\); retained fixed points belong to
\(\sigma\). The retained cycles and each choice are uniquely determined.
This proves the claimed bijection.

For \eqref{eq:roots}, the number of transpositions is \(a-|J|\).
Thus \(\sgn(\sigma)=\sgn(w)(-1)^{a-|J|}\), and this term contributes
\[
\sgn(w)(-1)^{a-|J|}
x^{b+|J|+|H|}u^{f-|H|}.
\]
Summing independently over \(J\subseteq E\) and \(H\subseteq F\) yields
\[
\sgn(w)x^b
\left(\sum_{J\subseteq E}(-1)^{a-|J|}x^{|J|}\right)
\left(\sum_{H\subseteq F}x^{|H|}u^{f-|H|}\right)
=\sgn(w)x^b(x-1)^a(u+x)^f.
\]
\end{proof}

\begin{proof}[Proof of the theorem]
The cases \(n=0,1\) are immediate, so use the preceding construction
for \(n\geq2\), first over \(\mathbb C\).
The sign map \(\varepsilon\) is an involutive algebra automorphism.
By Lemma~\ref{lem:coefficient},
\[
\varepsilon(\widehat S_n(u,x))=
\sum_{w\in S_n}x^{\anxc(w)}(x-1)^{\exc(w)}
(u+x)^{\fix(w)}w.
\]
In independent variables \(s,t\), substitute
\[
x=\frac{t}{t-s},\qquad u=\frac{1-t}{t-s}.
\]
Since \(x-1=s/(t-s)\), \(u+x=1/(t-s)\), and
\(\exc(w)+\anxc(w)+\fix(w)=n\), we obtain
\begin{equation}\label{eq:recover}
\sum_{a,b}s^at^bX_{a,b}
=(t-s)^n\varepsilon\!\left(
\widehat S_n\!\left(\frac{1-t}{t-s},\frac{t}{t-s}\right)\right).
\end{equation}
There is no residual denominator: a term \(C_{k,m}x^ku^m\) of
\(\widehat S_n\) becomes
\[
\varepsilon(C_{k,m})t^k(1-t)^m(t-s)^{n-k-m},
\]
where \(n-k-m\geq0\). Every coefficient on the right of
\eqref{eq:recover} therefore lies in the commutative algebra
\(\varepsilon(\A)\). Comparing coefficients gives
\(X_{a,b}\in\varepsilon(\A)\), proving commutativity over \(\mathbb C\).

For the asserted equality of algebras, put
\(B=\mathbb C[X_{a,b}:a,b\geq0]\).
The coefficient formula also expresses every coefficient of
\(\widehat S_n\) as a linear combination of the \(\varepsilon(X_{a,b})\):
group together permutations with the same \((a,b)\), for which
\(f=n-a-b\). Since those coefficients generate \(\A\), it follows that
\(\A\subseteq\varepsilon(B)\).
Together with \(B\subseteq\varepsilon(\A)\), this proves equality.

Finally each coefficient of \([X_{a,b},X_{c,d}]\) is an integer.
Its vanishing in \(\mathbb C[S_n]\) implies vanishing in
\(\mathbb Z[S_n]\). Base change proves the desired identity over
every commutative ring, including positive characteristic.
\end{proof}

\paragraph{Task 2. Construct a factorization bijection.}
For a fixed \(w\), construct an explicit bijection between
factorizations \(w=uv\) and \(w=v'u'\) preserving the respective
excedance--antiexcedance pairs of \(u,u'\) and \(v,v'\).

\paragraph{Task 3. Describe the joint spectrum combinatorially.}
Use the identified Bethe algebra to express the simultaneous eigenvalues
of the \(X_{a,b}\) in terms of explicit combinatorial data on each
irreducible symmetric-group representation.

\paragraph{Task 4. Determine the integral structure.}
Find generators and relations for the integral algebra
\(\mathbb Z[X_{a,b}]\), and determine its ranks, radicals and
semisimple quotients after reduction modulo primes.

\paragraph{Task 5. Find analogues for other groups and deformations.}
Determine which excedance statistics on signed permutations or which
deformations in Hecke algebras give commuting coefficient families,
and identify the corresponding transfer-polynomial formulas.

\begin{thebibliography}{9}
\bibitem{grinberg}
D. Grinberg, \emph{An introduction to the symmetric group algebra},
version 12 June 2026, Conjecture~3.10.22, p.~106,
\href{https://www.cip.ifi.lmu.de/~grinberg/t/24s/sga.pdf}{author's current text};
\href{https://arxiv.org/abs/2507.20706}{arXiv:2507.20706}.
\bibitem{mtv}
E. Mukhin, V. Tarasov and A. Varchenko,
\emph{Bethe subalgebras of the group algebra of the symmetric group},
Transform. Groups \textbf{18} (2013), 767--801,
\href{https://doi.org/10.1007/s00031-013-9232-y}{doi:10.1007/s00031-013-9232-y};
\href{https://arxiv.org/abs/1004.4248v4}{arXiv:1004.4248v4}.
\end{thebibliography}
\end{document}
